\documentclass{report}
\usepackage{amsmath,amssymb}
\setlength{\parindent}{0mm}
\setlength{\parskip}{1em}
\begin{document}
\begin{center}
\rule{6in}{1pt} \
{\large Pieterjan Robbe \\
{\bf A Practical Multilevel Quasi-Monte Carlo Method for PDEs with Random Coefficients}}

Department of Computer Science \\ KU Leuven - University of Leuven \\ Celestijnenlaan 200A \\ B3001 Leuven \\ Belgium
\\
{\tt pieterjan.robbe@kuleuven.be}\\
Dirk Nuyens\\
Stefan Vandewalle\end{center}

We present a Multilevel Quasi-Monte Carlo algorithm for the solution of
elliptic partial differential equations with random coefficients and
inputs. By combining the multilevel sampling idea with randomly shifted
rank-1 lattice rules, the algorithm constructs an estimator for the
expected value of some functional of the solution. The error analysis of
this estimator provides a formula for the optimal number of samples at
each level. The efficiency of the method is illustrated on a
three-dimensional subsurface flow problem with lognormal diffusion
coefficient. This example is particularly challenging because of the
small correlation length, and thus the large number of uncertainties that
must be included. We numerically show that it is possible to achieve a
cost almost inversely proportional to the requested tolerance on the
root-mean-square error.


\end{document}
